\section{CONCLUSÕES E TRABALHOS FUTUROS}

\subsection{Conclusões}

Este trabalho investigou a viabilidade de modelos de \textit{Gradient Boosting} --- XGBoost, LightGBM e CatBoost --- em um sistema de Aprendizado Federado para detecção de fraude bancária com restrições de equidade demográfica. O problema é intrinsecamente motivado por cenários reais onde dados de transações financeiras não podem ser compartilhados entre instituições por razões regulatórias e de privacidade, mas onde o treinamento colaborativo seria desejável para aumentar o poder preditivo. O dataset BAF \cite{jesus2022baf}, com aproximadamente um milhão de transações, taxa de fraude de 1\%, variável de grupo demográfico (faixa etária) e organização temporal, forneceu um \textit{benchmark} realista para a avaliação.

O resultado central é inequívoco: \textbf{a estratégia Cíclica superou consistentemente a estratégia Bagging} em desempenho (9--10 p.p. de TPR), estabilidade de convergência e eficiência de comunicação, em todos os três algoritmos avaliados. Esta superioridade, mais pronunciada com 50 rodadas do que em experimentos mais curtos reportados na literatura, tem origem no mecanismo de \textit{warm start} combinado com a agregação por ensemble da Bagging: ao acumular árvores progressivamente e selecionar apenas o modelo do melhor cliente por rodada, a Bagging induz \textit{over-boosting} que degrada o modelo com o passar das rodadas --- padrão \textit{crescimento--pico--degradação} claramente visível na Figura~\ref{fig:convergencia_tpr}. A estratégia Cíclica, ao contrário, distribui as atualizações sequencialmente entre clientes, produzindo um modelo coeso e estável que sustenta seu desempenho por todas as 50 rodadas avaliadas.

Entre os algoritmos, o LightGBM Cíclico obteve o melhor TPR fair absoluto (53,50\%), o XGBoost Cíclico a melhor fairness (FPR Ratio = 0,999 no regime fair, indicando equiparação virtual das taxas de falso positivo entre jovens e idosos --- posicionando-se no extremo direito do espaço fairness--performance da Figura~\ref{fig:tradeoff}), e o CatBoost Cíclico o melhor perfil de eficiência (104 segundos, 129 MB transferidos, conforme Figura~\ref{fig:eficiencia_tempo}) com desempenho competitivo (53,29\% TPR fair, FPR Ratio 0,972). Nos três casos, o regime fair --- que calibra limiares de decisão independentes por grupo demográfico --- elevou o FPR Ratio de 0,28--0,29 para 0,97--1,00 com perda de apenas 1,7--2,5 p.p. em TPR, demonstrando que equidade e desempenho são conciliáveis neste cenário com custo moderado. Os modelos Cíclicos situaram-se na faixa superior do \textit{baseline} centralizado do BAF (40--55\%), confirmando que o Aprendizado Federado com boosting não implica degradação severa de desempenho mesmo com dados fragmentados em clientes isolados.

\subsection{Contribuições}

As contribuições deste trabalho podem ser organizadas em três planos. No plano \textbf{empírico}, este trabalho fornece a primeira avaliação comparativa sistemática de XGBoost, LightGBM e CatBoost em ambiente federado com \textit{warm start} e restrições de fairness demográfica no dataset BAF, incluindo o comportamento de longo prazo das estratégias Bagging e Cíclica ao longo de 50 rodadas --- um horizonte temporal superior ao tipicamente avaliado na literatura de FL com boosting. No plano \textbf{metodológico}, foi demonstrada a compatibilidade entre otimização bayesiana federada de hiperparâmetros (Optuna com agregação por mediana), treinamento federado com \textit{warm start} e correção de fairness por \textit{post-processing}, compondo um \textit{pipeline} completo e reprodutível que vai da preparação dos dados à avaliação com equidade. No plano \textbf{prático}, os resultados fornecem diretrizes concretas para escolha de algoritmo e estratégia em sistemas de detecção de fraude com restrições de privacidade: a recomendação de CatBoost Cíclico com regime fair como configuração padrão, e a contraindicação da estratégia Bagging com muitas rodadas, são achados diretamente acionáveis por equipes de engenharia.

\subsection{Trabalhos Futuros}

Os resultados deste trabalho abrem várias direções de investigação. A limitação mais imediata é a avaliação exclusiva em cenário IID: experimentos com particionamento Non-IID --- simulando instituições com perfis distintos de clientes ou volumes de fraude desiguais --- são fundamentais para avaliar a robustez das conclusões em condições mais realistas. A literatura indica que a estratégia Cíclica tende a ser mais sensível à heterogeneidade dos dados do que a Bagging \cite{zhu2021noniid}, e verificar se isso anula a vantagem observada é uma pergunta direta de trabalho futuro.

Uma segunda direção é a incorporação de mecanismos de privacidade ao \textit{pipeline}: Privacidade Diferencial \cite{geyer2017differential} durante o treinamento ou Agregação Segura \cite{bonawitz2017secagg} durante a comunicação. Estes mecanismos são frequentemente necessários em produção para cumprir requisitos legais como o GDPR \cite{gdpr2016} e a LGPD, mas introduzem ruído que pode degradar o desempenho --- quantificar esse custo é essencial. Na mesma linha, a expansão para um número maior de clientes (10, 50, 100 instituições) permitiria avaliar a escalabilidade das estratégias: espera-se que a estratégia Cíclica aumente linearmente em tempo com o número de clientes, enquanto a Bagging mantém latência aproximadamente constante.

No plano algorítmico, há espaço para explorar variantes que mitiguem a degradação da Bagging, como o descarte periódico de árvores antigas (\textit{pruning}) ou a limitação do número máximo de estimadores transmitidos. Finalmente, a integração com infraestrutura de SDN (\textit{Software-Defined Networking}) --- motivação original desta linha de pesquisa --- permitiria avaliar se o controle programático da rede reduz a latência de comunicação percebida pelos clientes e se estratégias de priorização de tráfego alteram o comportamento de convergência dos experimentos aqui descritos.
